% ============================================================
% Chapter 5: Analysis
% ============================================================
\chapter{Analysis}
\label{ch:analysis}

\section{Statistical Analysis}
\label{sec:statistical-analysis}

All experimental results are reported as mean $\pm$ standard deviation over 20
independent runs.  We use the Wilcoxon signed-rank test (non-parametric,
paired) to assess statistical significance at the $\alpha = 0.05$ level, with
Bonferroni correction for multiple comparisons.  Effect sizes are reported
using Cohen's $d$.

\Cref{tab:statistical-tests} summarises the statistical comparison between
ASBO and each baseline across all benchmarks at $d=100$.

\begin{table}[htbp]
    \centering
    \caption{Statistical comparison of ASBO vs.\ baselines at $d=100$
    ($n=500$ evaluations).  $p$-values from Wilcoxon signed-rank test with
    Bonferroni correction.  ``+'' indicates ASBO is significantly better.}
    \label{tab:statistical-tests}
    \begin{tabularx}{\textwidth}{@{}lXXcX@{}}
        \toprule
        \textbf{Baseline} & \textbf{$p$-value} & \textbf{Cohen's $d$} &
        \textbf{Significant?} & \textbf{Improvement} \\
        \midrule
        REMBO    & $2.1 \times 10^{-6}$ & $1.82$ & Yes (+) & $3.2\times$ \\
        TuRBO    & $4.7 \times 10^{-7}$ & $2.14$ & Yes (+) & $2.8\times$ \\
        PESBO    & $1.3 \times 10^{-4}$ & $1.21$ & Yes (+) & $1.9\times$ \\
        SAASBO   & $8.9 \times 10^{-6}$ & $1.56$ & Yes (+) & $2.4\times$ \\
        \bottomrule
    \end{tabularx}
\end{table}

All comparisons are statistically significant ($p < 0.05$ after Bonferroni
correction) with large effect sizes ($d > 0.8$), confirming that the
improvements offered by ASBO are both statistically significant and practically
meaningful.

\section{Confidence Interval Analysis}
\label{sec:confidence-intervals}

We construct 95\% confidence intervals for the mean simple regret using the
bootstrap method (10{,}000 resamples).  \Cref{fig:confidence-intervals} shows
the confidence intervals for ASBO and baselines on the Ackley function at
$d=100$.

\begin{figure}[htbp]
    \centering
    \begin{tikzpicture}
        \begin{axis}[
            width=0.9\textwidth, height=5.5cm,
            xlabel={Method},
            ylabel={Simple regret (log scale)},
            ymode=log,
            xtick={1,2,3,4,5},
            xticklabels={ASBO, REMBO, TuRBO, PESBO, SAASBO},
            xmin=0.5, xmax=5.5,
            ymin=1e-4, ymax=1e-1,
            grid=major, grid style={gray!30},
            error bars/.cd, y dir=both, y explicit,
        ]
            \addplot+[only marks, mark=*, mark size=3pt, primary,
                error bars/.cd, y dir=both, y explicit]
                coordinates {
                    (1, 2.8e-3) +- (0, 3.2e-4) +- (0, 3.5e-4)
                    (2, 9.1e-3) +- (0, 1.1e-3) +- (0, 1.3e-3)
                    (3, 1.5e-2) +- (0, 2.8e-3) +- (0, 3.1e-3)
                    (4, 5.4e-3) +- (0, 7.2e-4) +- (0, 7.8e-4)
                    (5, 1.2e-2) +- (0, 1.8e-3) +- (0, 2.0e-3)
                };
        \end{axis}
    \end{tikzpicture}
    \caption{95\% confidence intervals for mean simple regret on Ackley
    ($d=100$, $n=500$).  Bootstrap with 10{,}000 resamples.}
    \label{fig:confidence-intervals}
\end{figure}

\section{Failure Mode Analysis}
\label{sec:failure-modes}

We identify three failure modes where ASBO's performance degrades:

\begin{enumerate}
    \item \textbf{High-frequency oscillations:} When the objective function
    exhibits rapid oscillations along the subspace directions, the \gls{gp}
    surrogate struggles to capture the local structure.  This is observed on
    the Rastrigin function for $d > 200$.

    \item \textbf{Subspace drift:} When the intrinsic dimensionality changes
    rapidly across the domain, the adaptive subspace may lag behind, leading
    to temporary misalignment.  This is mitigated by increasing the subspace
    update frequency.

    \item \textbf{Cold-start problem:} In the initial phase ($n < 30$), the
    subspace estimate is unreliable, and ASBO may waste evaluations exploring
    irrelevant dimensions.  This can be partially addressed by increasing the
    initial sample budget $n_0$.
\end{enumerate}

\section{Comparison with State of the Art}
\label{sec:sota-comparison}

\Cref{tab:sota-comparison} provides a comprehensive comparison of ASBO with
state-of-the-art methods across all 12 benchmarks at $d=100$, reporting the
median rank and the number of wins (best performance).

\begin{table}[htbp]
    \centering
    \caption{Comprehensive comparison at $d=100$ ($n=500$).  Median rank and
    wins across 12 benchmarks.}
    \label{tab:sota-comparison}
    \begin{tabularx}{\textwidth}{@{}Xccc@{}}
        \toprule
        \textbf{Method} & \textbf{Median Rank} & \textbf{Wins} &
        \textbf{Avg.\ Improvement} \\
        \midrule
        ASBO    & $\mathbf{1.0}$ & $\mathbf{10/12}$ & --- \\
        PESBO   & $2.5$ & $1/12$ & $1.9\times$ slower \\
        REMBO   & $3.0$ & $1/12$ & $3.2\times$ slower \\
        SAASBO  & $4.0$ & $0/12$ & $2.4\times$ slower \\
        TuRBO   & $4.5$ & $0/12$ & $2.8\times$ slower \\
        \bottomrule
    \end{tabularx}
\end{table}

\section{Case Studies}
\label{sec:case-studies}

\subsection{Molecular Property Optimisation}
\label{subsec:molecular}

In the molecular property optimisation task, ASBO discovers a drug candidate
with a predicted binding affinity of $-9.8$\,kcal/mol, compared to $-8.5$\,kcal/mol
for the best baseline.  The molecular structure, shown in \cref{fig:molecular},
features a novel binding motif not present in the training data.

\begin{figure}[htbp]
    \centering
    \begin{tikzpicture}
        \node[draw, rounded corners, fill=primary!5, minimum width=8cm,
              minimum height=3cm, align=center] {
            \textbf{Optimised Molecule}\\[6pt]
            Predicted binding affinity: $-9.8$\,kcal/mol\\
            Molecular weight: 423.5\,g/mol\\
            LogP: 2.8\\
            Drug-likeness (QED): 0.72
        };
    \end{tikzpicture}
    \caption{Summary of the best molecular candidate found by ASBO.}
    \label{fig:molecular}
\end{figure}

\subsection{Aerodynamic Shape Optimisation}
\label{subsec:aerodynamic}

For the airfoil drag minimisation, ASBO reduces the drag coefficient by 8.3\%
compared to the baseline NACA 0012 profile, using only 200 CFD evaluations
(where each evaluation requires approximately 4 hours of wall-clock time on a
high-performance computing cluster).
